\subsection{Особенности переходных процессов в области устойчивости и на ее границе}

Определим характер поведения системы при различных значениях коэффициентов регулятора $k_{\text{п}}$ и $k_{\text{и}}$ внутри области устойчивости и вблизи ее границ.

Рассмотрим сначала поведение системы при значениях
\[
k_{\text{п}} = 0.05, \qquad k_{\text{и}} = 0.05.
\]

Полученные графики переходных процессов представлены на рисунке~13.

\screenshot[0.75]{phi_pr_002_002}{Переходный процесс $\varphi_{\text{рпр}}(t)$ при $k_{\text{п}}=0.02$, $k_{\text{и}}=0.02$}

По полученным графикам видно, что система является устойчивой, однако переходный процесс протекает сравнительно медленно. Наблюдается небольшое перерегулирование, после чего колебания быстро затухают. Время переходного процесса составляет порядка $50$~с.

Рассмотрим систему при значениях коэффициентов
\[
k_{\text{п}} = 5, \qquad k_{\text{и}} = 10.
\]

Полученные графики представлены на рисунке~14.

\screenshot[0.8]{process_5_10}{Переходные процессы при $k_{\text{п}} = 5$, $k_{\text{и}} = 10$}

При данных коэффициентах переходный процесс становится существенно более колебательным. Амплитуда колебаний уменьшается со временем, однако затухание происходит медленно. В системе возникают значительные динамические нагрузки, а длительность переходного процесса составляет приблизительно $283$~с.

Рассмотрим поведение системы на другой границе области устойчивости. Примем
\[
k_{\text{п}} = 5, \qquad k_{\text{и}} = 75.
\]

Соответствующие графики приведены на рисунке~15.

\screenshot[0.8]{process_5_75}{Переходные процессы при $k_{\text{п}} = 5$, $k_{\text{и}} = 75$}

Из полученных результатов видно, что система остается устойчивой, однако колебательный характер переходного процесса выражен значительно сильнее. Амплитуда колебаний уменьшается крайне медленно, вследствие чего длительность переходного процесса возрастает примерно до $1410$~с. Данный режим находится вблизи границы устойчивости и характеризуется низким запасом устойчивости.
